考虑一维定常、平面、绝热的理想气体法向激波。所有速度均相对于静止激波面测量；远离薄激波层的两侧均为均匀平衡态，忽略黏性应力与热流在两侧控制面的通量，但允许激波内部不可逆耗散。上游量为 \(p_1>0,\rho_1>0,u_1>0\)，绝热指数 \(\gamma>1\) 为常数，比气体常数 \(R>0\)，定压比热 \(c_p=\gamma R/(\gamma-1)\)。上游声速 \(c_1=\sqrt{\gamma p_1/\rho_1}\)，马赫数 \(M_1=u_1/c_1>1\)。下游用下标 \(2\) 表示，定义 \(x=M_1^2\)、密度比 \(r=\rho_2/\rho_1\)、压强比 \(P=p_2/p_1\)。
(1) 从控制体的质量、动量和总焓通量守恒出发，推导质量通量 \(j\)、Rayleigh 动量关系和 Hugoniot 能量关系，并说明为何不能额外强加等熵条件。
(2) 独立消元求 \(r,P,T_2/T_1\) 和下游马赫数 \(M_2\) 的闭式表达式；验证 \(M_1>1\) 的压缩激波必有 \(r>1,P>1,0<M_2<1\)。
(3) 从理想气体熵函数计算每单位质量的熵增 \(\Delta s\)，严格证明其为正；在弱激波极限 \(x=1+\varepsilon\)、\(0<\varepsilon\ll1\) 下，求 \(\Delta s/R\) 的首个非零项及其阶数。
(4) 求 \(M_1\to\infty\) 时的最大压缩比、压强和温度的主导渐近、下游马赫数极限；分别用守恒方程和量纲分析核对结果，并解释强激波中动能转为内能的物理机制。
